% ===== §12 A Typology of Everyday Aesthetic Practice =====
\section{A Typology of Everyday Aesthetic Practice}
\label{p23:sec:typology}

\noindent
This section makes the paper's most concrete offering: four models of everyday aesthetic practice, proposed to occupy the space the survey left open, generative rather than reproductive or resistant, relational rather than individual, conducted in the flow of ordinary shared life, and stated in the general terms of the theory the paper has built. Each is given with its theoretical descent made explicit, because a model of practice that could not say what it rests on would be a suggestion rather than a proposal, and the paper's claim is that these four are derived from its theory and not improvised. The models are offered as preliminary and open, a beginning of the typology rather than its closure, and the reason for their openness is internal to the fourth of them, which makes the openness of an aesthetic practice a structural requirement rather than a modesty. The four share a single skeleton, and the section states it first, then develops each model in turn with its basis and its ethics, and closes on the unity of the typology and the sense in which it remains unfinished.

\subsection{The common skeleton}
\label{p23:sub:typ-skeleton}

The four models are variations on one structure, and seeing the structure first prevents the models from looking like a list. Each moves from a perception or its equivalent, through a reception, through a sharing that enters the relational system, to a reproduction or feedback that returns upon the shared life and begins the cycle again. The skeleton is the good cycle of value in the domain of practice: a generative circuit that returns to its root with a surplus around an unsettled centre, whose positive holonomy is the mark of its goodness (\xref{p23:sub:sym-criterion}). What distinguishes the models is where they enter the circuit and what happens at each of its moments. The first enters through active perception and is the paradigm. The second enters through silence and makes reception the suspension of judgement. The third does not directly produce at all, but constructs the field in which the circuit may generate of itself. The fourth iterates the circuit across many turns, refining the strategy by which it is entered. Three of the four, the first, second, and fourth, are the same circuit run at different timescales, single, reflective, and evolutionary; the third is the circuit's complement, the making of room for it. Together they cover the principal temporalities of everyday aesthetic generation, and none is advanced as excluding the others, since a shared life will move among them.

\subsection{Model I: perceiving, receiving, sharing, reproducing}
\label{p23:sub:typ-model1}

\begin{claim}[Model I]
The basic generative cycle of everyday aesthetic practice runs from a shared perception, through its reception into the perceivers, through its symbolisation and sharing, to its reproduction in the relation, which returns the perceivers to a deepened perception and turns the cycle again.
\end{claim}

The theoretical basis of the first model is threefold and has been established in earlier parts of the paper. Its rhythm is Dewey's: the movement of doing and undergoing by which an experience is brought to fullness, taken here into the shared case, so that the perceiving and the receiving are the undergoing and the sharing and reproducing are the doing of a common aesthetic experience \parencite{dewey1934}. Its value structure is the good cycle of this series: the circuit generates a surplus around an unsettled centre and returns it to the shared attending, a positive holonomy in the register of practice \parencite{paperix,berry1984}. And its ethics is the theory of relational reproduction: the reproduction that closes the cycle must satisfy the ordered principles of that theory, generativity, co-reproduction, non-exploitation, counter-power, and the preservation of the subject, so that what is reproduced is a shared perception in which each remains a generating subject and neither is reduced to the material of the other's experience \parencite{wanreproduction}. The model has, accordingly, two points of ethical exposure, and naming them is part of specifying it. The sharing carries an ethics of sharing: to share a perception is to symbolise it toward another, and the sharing is generative only if it is an invitation to perceive together rather than a display that makes the other the audience of one's own seeing, which is the aesthetic form of the epistemic hospitality this series has developed, the translation of a perception into the other's sensibility as a respect for it rather than a management of it. The reproduction carries an ethics of reproduction: the perception reproduced in the relation must leave the other room to generate, not settle into a reproduced taste imposed on them, for a reproduction that fixes the other's perception has closed the cycle it was meant to turn. The first model is the paradigm because it is the full circuit run once, with both its generative promise and its two ethical exposures in plain view.

\subsection{Model II: silence, receiving, sharing, relational reproduction}
\label{p23:sub:typ-model2}

\begin{claim}[Model II]
A reflective variant of the cycle begins in individual silence, which is itself a reproduction of experience; its reception is the suspension of private judgement; its sharing returns the unsettled experience to the relation; and the relational reproduction completes the judgement in common rather than in the private interior.
\end{claim}

The second model shares the skeleton of the first and differs at two moments, and its theoretical basis is what makes the difference intelligible. It begins not in active perception but in individual silence, and its first claim is that this silence is not an absence of activity but a reproduction of experience in its own right. This is the psychoanalytic contribution the paper reserved: the working-through by which an experience is reworked after the fact, and the deferred action, \emph{Nachträglichkeit}, by which an experience receives its sense belatedly, in a later revisiting that reproduces it as something it was not at first \parencite{lacan1978}. Silence, in this model, is the interval in which this internal reproduction occurs. Its danger was marked earlier and is met here by the model's second distinctive moment. What silence reproduces may be understanding and gratitude or may be resentment and misreading, and a silence that hardened its rumination into a private verdict would sediment the bad reproduction before it ever reached the relation, which is the mechanism of the cold estrangement that this series has named among the insidious injuries of a shared life. The model's reception is the guard against this. Reception here is the suspension of private judgement, the epoch\'e in which the perceiver holds what the silence has surfaced without ruling on its value in the private interior, and carries it, unsettled, back to the relation to be judged in common \parencite{husserl2001}. This is the aesthetic instance of the presentational justice this series developed for the non-binding vow: the individual does not sit as the sole judge of their own experience in a private court, but suspends the verdict and returns the matter to the relational system in which alone it is rightly adjudicated. The ethics of the second model is thus built into its structure rather than added to it, for the reception that suspends private judgement is the very device by which the silent reproduction is kept from going bad; the good silence and the bad are told apart by whether the silence passes through this suspension or hardens into a private verdict. And the model demonstrates something about the skeleton it shares with the first, which is that the cycle does not depend on an external aesthetic stimulus to drive it: the reworking of an experience already had, in silence, can enter and turn the same circuit, so that the generativity of the shared cycle is shown to be capable of running on the internal reproduction of experience alone.

\subsection{Model III: the leaving of emptiness}
\label{p23:sub:typ-model3}

\begin{claim}[Model III]
Rather than producing a perception, one may construct the field in which perception generates of itself, leaving an emptiness in which the contingency of the everyday is given room to occur.
\end{claim}

The third model is the complement of the others, and its theoretical basis is the account of the field of contingency joined to the Daoist ethics of non-possession. Where the first two models enter and turn the generative circuit, the third makes room for it, and this making of room has been prepared twice in the paper: in the diagnosis, where the everyday was shown to be the field of contingency, generative precisely because it is not arranged in advance (\xref{p23:sub:diag-contingency}), and in the design of space, where the interval, the \zh{間}, was shown to be a positive element of composition, the built form of that contingency (\xref{p23:sub:space-atmosphere}). To leave emptiness is to construct, in the arrangement of a shared life and its spaces, the room in which perception may generate of itself rather than being directed to a prescribed object. Its ethics is the Daoist \zh{玄德}, the dark virtue that generates without possessing, acts without presuming, and leads without commanding, \zh{生而不有，為而不恃，長而不宰} \parencite{laozi1963}. The leaving of emptiness is the practical form of this virtue in everyday aesthetic life: the restraint that declines to produce the perception, that does not fill the shared field with one's own arrangement, that trusts the contingency of the ordinary and the other's generativity to bring forth what no one commanded. Its characteristic danger is the contrary of the second model's: not the hardening of a private verdict but the abdication that mistakes mere absence for generative emptiness, leaving not a field but a void. The virtue is not in the emptiness as such but in the emptiness constructed as a field, tuned like B\"{o}hme's atmosphere and composed like the interval of \emph{ma}, so that it affords generation rather than merely withholding presence. The leaving of emptiness is the hardest of the models to do well, because it must produce, through restraint, the conditions of a generation it does not itself perform.

\subsection{Model IV: iterative symbolic refinement}
\label{p23:sub:typ-model4}

\begin{claim}[Model IV]
Because the good of a symbolic arrangement cannot be known in advance, everyday aesthetic practice may proceed by iterated refinement: an aesthetic perception issues in a symbolic strategy, the strategy is enacted in a symbolic arrangement, the arrangement enters the relational system, the relation returns a value-feedback, and the strategy is refined for the next turn.
\end{claim}

The fourth model is the paper's most distinctive proposal, and it exists to solve a circularity the other models leave standing. To arrange the symbols of a shared life well, a love letter, a note left in the morning, the small benevolent signs by which a home is made to generate tenderness, would seem to require that one already know what a good such arrangement is; but the good aesthetic arrangement is exactly what an aesthetic education has yet to generate, and cannot be presupposed without begging the question the whole paper is about. The fourth model dissolves this circularity by making operable not the good, which cannot be operated, but the symbolic strategy, which can. One does not produce a good sign; one enacts a symbolic strategy, returns it to the relation, receives the relation's response, and refines the strategy accordingly. The theoretical basis of this is evolutionary, and it is precise rather than metaphorical. The structure is that of an evolutionary process, variation of strategy, selection, and refinement across iterations; but it is an evolutionary process of a particular kind, one without a fitness function, because the aesthetic has no objective function that could score a strategy in advance. Here the model draws on the finding of Lehman and Stanley that in a search whose objective cannot be well specified, and whose direct pursuit is therefore deceptive, abandoning the objective and searching instead by other signals can reach what the pursuit of the objective could not \parencite{lehman2011}. The aesthetic is such a search: the direct pursuit of beauty, the arrangement made to be beautiful, is deceptive in exactly their sense, as the photographed and displayed meal is deceived out of the beauty it grasped at. The fourth model replaces the missing fitness function not with novelty as such but with the value-feedback of the relation, which is not a scalar score but a multidimensional and generative response, and which is itself in evolution, since what counts as a good response is redefined by the relation as it develops. The good, in this model, is not presupposed and then approximated; it emerges through the co-evolution of the strategy and the relation's own criterion for it, so that the aesthetic education of a shared symbolic life is an open-ended evolution without a fixed fitness landscape. Its ethics is the ethics of the symbolic strategy, and it descends from the symbolic order's two fates (\xref{p23:sec:symbolic}): the refinement is generative where the value-feedback returns to both parties and appropriative where the strategy is refined toward the management or manipulation of the other, so that the direction of refinement, toward the positive holonomy of a shared surplus or the negative holonomy of a one-sided capture, is what tells the good iteration from the bad. This model, finally, is the reason the typology is offered as open. If the good aesthetic arrangement is not presupposed but emerges through open-ended refinement, then a fixed and closed set of models would contradict the very structure the fourth model describes, and the typology must itself remain, like the practice it theorises, open to the refinement of its own strategies.

\subsection{The unity of the typology}
\label{p23:sub:typ-unity}

The four models compose one typology, and its unity can now be stated against the limits the survey fixed. They cover four temporalities of everyday aesthetic generation: the active generation of the first, the generation after reflective sedimentation of the second, the making of room for generation of the third, and the generation across iterated refinement of the fourth. The first, second, and fourth are the shared circuit run at the timescales of the single occasion, the reworked experience, and the evolutionary series; the third is the circuit's complement, the constructed emptiness in which it may occur. Each passes a limit the survey named: against the theory of practice, which gave everyday doing as reproduction or resistance, all four are generative, producing new aesthetic value rather than reproducing a structure or resisting one; against the methods of aesthetic education, which trained the individual in a reserved space, all four work the relation in the flow of ordinary life; against the sub-field of everyday aesthetics, which treated the everyday aesthetic as an individual's private experience, all four are relational, taking the shared generation of perception as what they cultivate; and against the practice traditions of the East, which enacted such generation within particular frameworks, all four are stated in the general terms of the paper's theory of perception and value, as the structure those traditions have long enacted. The typology is preliminary, and it is preliminary in the strong sense the fourth model requires: not a finished inventory awaiting application but an opening set, offered for the refinement that its own fourth model makes the condition of any living aesthetic practice. What the models share, beneath their differences, is the single commitment from which the whole paper proceeds, that the unit of everyday aesthetic practice is the relation, and that to practise the aesthetics of the everyday is to turn, between two people, the generative cycle by which their shared and ordinary world is continually brought to perception.
